\section{Inner product, length, and orthogonality}
For real columns $u,v$ of equal length, $u\cdot v=u^Tv$ is the sum of
corresponding coordinate products. The norm is
$\|v\|=\sqrt{v\cdot v}$, and distance is $\|u-v\|$.
Two vectors are orthogonal precisely when their inner product is zero.
Normalize a nonzero vector by dividing by its length, not by its squared
length. The zero vector has no unit vector in its direction.

For a subspace $W$, its orthogonal complement $W^\perp$ contains the
vectors orthogonal to \emph{every} vector of $W$. Testing a spanning family
suffices because the inner product distributes over linear combinations.

\practice{6.1.P1}{A projection coefficient and its scaled vector}
Compute $a\cdot b=(-2)(-3)+1(1)=7$ and
$a\cdot a=(-2)^2+1^2=5$. Thus
\[
\frac{a\cdot b}{a\cdot a}=\frac75,\qquad
\left(\frac{a\cdot b}{a\cdot a}\right)a
=\begin{bmatrix}-14/5\\7/5\end{bmatrix}.
\]
The first answer is a scalar, while the second is a vector.

\practice{6.1.P2}{Normalize the three-coordinate vector}
\step{Compute the length, keeping the square root nonnegative.}
\[\|v\|=\sqrt{(\frac{4}{3})^2+(-1)^2+(\frac{2}{3})^2}=\frac{\sqrt{29}}{3}.\]
Since $v\ne0$, divide by this positive length rather than by its square:
\[u=\frac{v}{\|v\|}=\left[\begin{array}{r}\frac{4 \sqrt{29}}{29}\\- \frac{3 \sqrt{29}}{29}\\\frac{2 \sqrt{29}}{29}\end{array}\right].\]
The multiplier $1/\|v\|$ is positive, so the direction is unchanged; also $u\cdot u=(v\cdot v)/\|v\|^2=1$.
\answer{The displayed $u$ is the unit vector in the given direction.}

\practice{6.1.P3}{Check orthogonality directly}
The inner product is
$5(4/3)+6(-1)+(-1)(2/3)=20/3-6-2/3=0$.
Therefore $d$ is orthogonal to $c$; checking the zero inner product is
exactly the definition needed.

\practice{6.1.P4}{Orthogonality survives normalization}
The unit vector $u$ from Practice 2 is $c/\|c\|$.
By linearity in the second argument,
$d\cdot u=(1/\|c\|)(d\cdot c)=0$ using Practice 3.
\answer{Multiplying one of two orthogonal vectors by a scalar does
not destroy orthogonality.}

\exercise{6.1.1}{Inner-product and norm calculations}
\step{Compute the two dot products from corresponding coordinates.}
\[u\cdot u=(-1)(-1)+(2)(2)=5,\]
\[v\cdot u=(4)(-1)+(6)(2)=8.\]
\answer{The requested ratio is $(v\cdot u)/(u\cdot u)=\frac{8}{5}$.}

\exercise{6.1.2}{Inner-product and norm calculations}
\step{Compute the two dot products from corresponding coordinates.}
\[w\cdot w=(3)(3)+(-1)(-1)+(-5)(-5)=35,\]
\[x\cdot w=(6)(3)+(-2)(-1)+(3)(-5)=5.\]
\answer{The requested ratio is $(x\cdot w)/(w\cdot w)=\frac{1}{7}$.}

\exercise{6.1.3}{Inner-product and norm calculations}
The squared length is $w\cdot w=35$. Multiply every coordinate by its reciprocal:
\[\frac1{w\cdot w}w=\frac{1}{35}\left[\begin{array}{r}3\\-1\\-5\end{array}\right]=\left[\begin{array}{r}\frac{3}{35}\\- \frac{1}{35}\\- \frac{1}{7}\end{array}\right].\]
This is not generally a unit vector: normalization would divide by the square root of this denominator.

\exercise{6.1.4}{Inner-product and norm calculations}
The squared length is $u\cdot u=5$. Multiply every coordinate by its reciprocal:
\[\frac1{u\cdot u}u=\frac{1}{5}\left[\begin{array}{r}-1\\2\end{array}\right]=\left[\begin{array}{r}- \frac{1}{5}\\\frac{2}{5}\end{array}\right].\]
This is not generally a unit vector: normalization would divide by the square root of this denominator.

\exercise{6.1.5}{Inner-product and norm calculations}
The numerator is $u\cdot v=8$ and the denominator is $v\cdot v=52$. Therefore
\[\left(\frac{u\cdot v}{v\cdot v}\right)v=\frac{2}{13}\left[\begin{array}{r}4\\6\end{array}\right]=\left[\begin{array}{r}\frac{8}{13}\\\frac{12}{13}\end{array}\right].\]
The denominator is nonzero because the vector being projected onto is nonzero.

\exercise{6.1.6}{Inner-product and norm calculations}
The numerator is $w\cdot x=5$ and the denominator is $x\cdot x=49$. Therefore
\[\left(\frac{w\cdot x}{x\cdot x}\right)x=\frac{5}{49}\left[\begin{array}{r}6\\-2\\3\end{array}\right]=\left[\begin{array}{r}\frac{30}{49}\\- \frac{10}{49}\\\frac{15}{49}\end{array}\right].\]
The denominator is nonzero because the vector being projected onto is nonzero.

\exercise{6.1.7}{Inner-product and norm calculations}
By the norm definition,
\[\|w\|=\sqrt{(3)^2+(-1)^2+(-5)^2}=\sqrt{35}=\sqrt{35}.\]
The norm is the nonnegative square root, so there is no $\pm$ choice.

\exercise{6.1.8}{Inner-product and norm calculations}
By the norm definition,
\[\|x\|=\sqrt{(6)^2+(-2)^2+(3)^2}=\sqrt{49}=7.\]
The norm is the nonnegative square root, so there is no $\pm$ choice.

\exercise{6.1.9}{Normalize without changing the direction}
\step{Compute the length, keeping the square root nonnegative.}
\[\|v\|=\sqrt{(-30)^2+(40)^2}=50.\]
Since $v\ne0$, divide by this positive length rather than by its square:
\[u=\frac{v}{\|v\|}=\left[\begin{array}{r}- \frac{3}{5}\\\frac{4}{5}\end{array}\right].\]
The multiplier $1/\|v\|$ is positive, so the direction is unchanged; also $u\cdot u=(v\cdot v)/\|v\|^2=1$.
\answer{The displayed $u$ is the unit vector in the given direction.}

\exercise{6.1.10}{Normalize without changing the direction}
\step{Compute the length, keeping the square root nonnegative.}
\[\|v\|=\sqrt{(-6)^2+(4)^2+(-3)^2}=\sqrt{61}.\]
Since $v\ne0$, divide by this positive length rather than by its square:
\[u=\frac{v}{\|v\|}=\left[\begin{array}{r}- \frac{6 \sqrt{61}}{61}\\\frac{4 \sqrt{61}}{61}\\- \frac{3 \sqrt{61}}{61}\end{array}\right].\]
The multiplier $1/\|v\|$ is positive, so the direction is unchanged; also $u\cdot u=(v\cdot v)/\|v\|^2=1$.
\answer{The displayed $u$ is the unit vector in the given direction.}

\exercise{6.1.11}{Normalize without changing the direction}
\step{Compute the length, keeping the square root nonnegative.}
\[\|v\|=\sqrt{(\frac{7}{4})^2+(\frac{1}{2})^2+(1)^2}=\frac{\sqrt{69}}{4}.\]
Since $v\ne0$, divide by this positive length rather than by its square:
\[u=\frac{v}{\|v\|}=\left[\begin{array}{r}\frac{7 \sqrt{69}}{69}\\\frac{2 \sqrt{69}}{69}\\\frac{4 \sqrt{69}}{69}\end{array}\right].\]
The multiplier $1/\|v\|$ is positive, so the direction is unchanged; also $u\cdot u=(v\cdot v)/\|v\|^2=1$.
\answer{The displayed $u$ is the unit vector in the given direction.}

\exercise{6.1.12}{Normalize without changing the direction}
\step{Compute the length, keeping the square root nonnegative.}
\[\|v\|=\sqrt{(\frac{8}{3})^2+(2)^2}=\frac{10}{3}.\]
Since $v\ne0$, divide by this positive length rather than by its square:
\[u=\frac{v}{\|v\|}=\left[\begin{array}{r}\frac{4}{5}\\\frac{3}{5}\end{array}\right].\]
The multiplier $1/\|v\|$ is positive, so the direction is unchanged; also $u\cdot u=(v\cdot v)/\|v\|^2=1$.
\answer{The displayed $u$ is the unit vector in the given direction.}

\exercise{6.1.13}{Inner-product and norm calculations}
Subtract the point coordinates first:
\[a-b=\left[\begin{array}{r}10\\-3\end{array}\right]-\left[\begin{array}{r}-1\\-5\end{array}\right]=\left[\begin{array}{r}11\\2\end{array}\right].\]
\answer{The distance is $\|a-b\|=\sqrt{125}=5 \sqrt{5}$. Reversing the subtraction gives the same length.}

\exercise{6.1.14}{Inner-product and norm calculations}
Subtract the point coordinates first:
\[a-b=\left[\begin{array}{r}0\\-5\\2\end{array}\right]-\left[\begin{array}{r}-4\\-1\\8\end{array}\right]=\left[\begin{array}{r}4\\-4\\-6\end{array}\right].\]
\answer{The distance is $\|a-b\|=\sqrt{68}=2 \sqrt{17}$. Reversing the subtraction gives the same length.}

\exercise{6.1.15}{Inner-product and norm calculations}
Compute the defining dot product:
\[a\cdot b=(8)(-2)+(-5)(-3)=-1.\]
\answer{The pair is not orthogonal because the product is nonzero.}

\exercise{6.1.16}{Inner-product and norm calculations}
Compute the defining dot product:
\[a\cdot b=(12)(2)+(3)(-3)+(-5)(3)=0.\]
\answer{The pair is orthogonal because the product is zero.}

\exercise{6.1.17}{Inner-product and norm calculations}
Compute the defining dot product:
\[a\cdot b=(3)(-4)+(2)(1)+(-5)(-2)+(0)(6)=0.\]
\answer{The pair is orthogonal because the product is zero.}

\exercise{6.1.18}{Inner-product and norm calculations}
Compute the defining dot product:
\[a\cdot b=(-3)(1)+(7)(-8)+(4)(15)+(0)(-7)=1.\]
\answer{The pair is not orthogonal because the product is nonzero.}

\exercise{6.1.19}{Inner-product and norm calculations}
\begin{enumerate}[label=(\alph*)]
\item True. It is the squared-norm identity $\|v\|^2=v\cdot v$.
\item True. Scalar distributivity gives $u\cdot(cv)=c(u\cdot v)$.
\item True. Squaring the equal distances and expanding gives
$\|u\|^2-2u\cdot v+\|v\|^2=\|u\|^2+2u\cdot v+\|v\|^2$.
Hence $4u\cdot v=0$, so the vectors are orthogonal.
\item False. Take $A=\begin{bmatrix}1&1\\0&0\end{bmatrix}$.
Its column space contains $(1,0)^T$ and its null space contains
$(-1,1)^T$; their dot product is $-1$, not zero. The correct general
pairing is \emph{row space} with null space.
\item True. For $w=\sum_jc_jv_j$ we have
$x\cdot w=\sum_jc_j(x\cdot v_j)=0$; thus $x\in W^\perp$.
\end{enumerate}

\exercise{6.1.20}{Inner-product and norm calculations}
\begin{enumerate}[label=(\alph*)]
\item True. Symmetry gives $u\cdot v=v\cdot u$, so their difference
is zero.
\item False for negative scalars. For $c=-1$ and $v\ne0$,
$\|cv\|=\|v\|$ but $c\|v\|=-\|v\|$. The correct rule uses $|c|$.
\item True. Being orthogonal to every member of $W$ is the definition
of membership in $W^\perp$.
\item True. Expand $\|u+v\|^2=\|u\|^2+2u\cdot v+\|v\|^2$.
The given equality makes $2u\cdot v=0$.
\item True. $Ax=0$ means every row has dot product zero with $x$.
The same holds for every linear combination of rows, hence all of Row $A$.
\end{enumerate}

\exercise{6.1.21}{Inner-product and norm calculations}
\step{(b) Distribute using the transpose and matrix multiplication rules.}
\[
(u+v)\cdot w=(u+v)^Tw=(u^T+v^T)w=u^Tw+v^Tw
=u\cdot w+v\cdot w.
\]
Here we used $(u+v)^T=u^T+v^T$ and distribution of matrix products
across addition, both from Chapter 2.
\step{(c) Move a scalar through the products.}
The transpose rule $(cu)^T=cu^T$ and scalar associativity of products give
\[
(cu)\cdot v=(cu)^Tv=c(u^Tv)=c(u\cdot v)
=u^T(cv)=u\cdot(cv).
\]
Each expression is a $1\times1$ product, identified with its real scalar value.

\exercise{6.1.22}{Inner-product and norm calculations}
We have $u\cdot u=u_1^2+u_2^2+u_3^2$. Every real square is
nonnegative, so the sum is nonnegative. If the sum is zero, none of the
nonnegative summands can be positive, so $u_i^2=0$ for every $i$.
This forces all three entries to be zero. Conversely the zero vector has
zero inner product with itself.
\answer{$u\cdot u\ge0$, with equality if and only if $u=0$.}

\exercise{6.1.23}{Inner-product and norm calculations}
Compute directly rather than invoking the Pythagorean theorem:
\[u\cdot v=(2)(-7)+(-5)(-4)+(-1)(6)=0,\]
\[\|u\|^2=30,\qquad\|v\|^2=101,\qquad u+v=\left[\begin{array}{r}-5\\-9\\5\end{array}\right].\]
\[\|u+v\|^2=(-5)^2+(-9)^2+(5)^2=131=30+101.\]
The dot product is zero, so this example agrees with the orthogonal-sum identity, which has now been verified by arithmetic.

\exercise{6.1.24}{Inner-product and norm calculations}
Use bilinearity and symmetry to expand both squared norms:
\[
\|u+v\|^2=\|u\|^2+2u\cdot v+\|v\|^2,
\]
\[
\|u-v\|^2=\|u\|^2-2u\cdot v+\|v\|^2.
\]
Adding cancels the two cross terms and gives
$\|u+v\|^2+\|u-v\|^2=2\|u\|^2+2\|v\|^2$.
\answer{The parallelogram law holds for all $u,v\in\R^n$;
orthogonality is not assumed.}

\exercise{6.1.25}{Inner-product and norm calculations}
Orthogonality to $(a,b)^T$ is the single homogeneous equation
$ax+by=0$. If $a=b=0$, every $(x,y)$ works, so $H=\R^2$.
If $(a,b)\ne(0,0)$, the nonzero vector $(-b,a)^T$ is orthogonal
to $(a,b)^T$ and its multiples solve the equation. Conversely, if
$a\ne0$ solve $x=-(b/a)y$; if $a=0$ then $b\ne0$ and the
equation gives $y=0$. Either way every solution is a multiple of
$(-b,a)^T$.
\answer{For $v\ne0$, $H$ is the line through the origin perpendicular
to $v$; for $v=0$, it is the entire plane.}

\exercise{6.1.26}{Inner-product and norm calculations}
We can write $W=\Nul[5\ {-6}\ 7]$. The Chapter 4 theorem that
a null space is a subspace therefore applies. Explicitly its equation is
$5x_1-6x_2+7x_3=0$, one nonzero homogeneous linear constraint.
Solving for $x_1$ leaves two freely chosen coordinates, so the geometric
set is a plane through the origin. The normal direction is $u$ because
$u\cdot x=0$ for every vector in the plane.
\answer{$W$ is the plane perpendicular to $(5,-6,7)^T$ through zero.}

\exercise{6.1.27}{Inner-product and norm calculations}
By distributivity,
$y\cdot(u+v)=y\cdot u+y\cdot v=0+0=0$.
The two zero values follow from the two given orthogonality assumptions.
\answer{$y$ is orthogonal to $u+v$.}

\exercise{6.1.28}{Inner-product and norm calculations}
Take any $w\in\Span\{u,v\}$. There are real scalars $c_1,c_2$
with $w=c_1u+c_2v$. Then
$y\cdot w=c_1(y\cdot u)+c_2(y\cdot v)=0$.
Because $w$ was arbitrary, the claim holds for every vector in the span,
not merely for its two generators.
\answer{$y\in(\Span\{u,v\})^\perp$.}

\exercise{6.1.29}{Inner-product and norm calculations}
An arbitrary $w\in W$ can be written $w=\sum_{j=1}^pc_jv_j$.
Apply distributivity repeatedly:
\[
x\cdot w=\sum_{j=1}^pc_j(x\cdot v_j)=\sum_{j=1}^pc_j0=0.
\]
Thus every $w$ in $W$ is orthogonal to $x$.
\answer{Testing orthogonality against a spanning family tests all of $W$.}

\exercise{6.1.30}{Inner-product and norm calculations}
\step{(a) Closure under scalar multiplication.}
Let $z\in W^\perp$, $c\in\R$, and let $u$ be any vector of $W$.
Then $(cz)\cdot u=c(z\cdot u)=c0=0$. Since this holds for every
$u\in W$, $cz\in W^\perp$.
\step{(b) Closure under addition.}
For $z_1,z_2\in W^\perp$ and any $u\in W$,
$(z_1+z_2)\cdot u=z_1\cdot u+z_2\cdot u=0+0=0$.
Hence $z_1+z_2\in W^\perp$.
\step{(c) Include zero and apply the subspace test.}
The zero vector has zero dot product with every $u\in W$, so it lies
in $W^\perp$. Together with (a) and (b), these are all three required
subspace properties.
\answer{$W^\perp$ is a subspace of $\R^n$.}

\exercise{6.1.31}{Inner-product and norm calculations}
If $x\in W$ and $x\in W^\perp$, then the latter membership makes
$x$ orthogonal to every member of $W$, including itself. Thus
$x\cdot x=\|x\|^2=0$, which forces $x=0$ by positive definiteness.
\answer{$W\cap W^\perp=\{0\}$.}

\exercise{6.1.32}{Inner-product and norm calculations}
\step{(a) Check all column lengths and mutual dot products.}
The single exact calculation $A^TA=I_4$ states that every column has squared length $1$ and all six off-diagonal pairwise products are zero. Thus the four lengths are $1,1,1,1$ and the six requested dot products are $0$.
\step{(b)--(d) Perform three reproducible random experiments.}
Using integer draws from $[-3,3]$ with Python random seed $61032$ gives the following three nonzero pairs. Listed equalities compare the original and transformed quantities.
\[u=\left[\begin{array}{r}2\\2\\0\\-3\end{array}\right],\quad v=\left[\begin{array}{r}-1\\-2\\1\\0\end{array}\right];\]
\[\|u\|=\|Au\|=\sqrt{17},\quad\|v\|=\|Av\|=\sqrt{6},\]
\[\cos\angle(u,v)=\cos\angle(Au,Av)=- \frac{\sqrt{102}}{17}.\]
\[u=\left[\begin{array}{r}-1\\0\\3\\2\end{array}\right],\quad v=\left[\begin{array}{r}3\\0\\2\\1\end{array}\right];\]
\[\|u\|=\|Au\|=\sqrt{14},\quad\|v\|=\|Av\|=\sqrt{14},\]
\[\cos\angle(u,v)=\cos\angle(Au,Av)=\frac{5}{14}.\]
\[u=\left[\begin{array}{r}-1\\1\\1\\-1\end{array}\right],\quad v=\left[\begin{array}{r}0\\2\\-1\\-1\end{array}\right];\]
\[\|u\|=\|Au\|=2,\quad\|v\|=\|Av\|=\sqrt{6},\]
\[\cos\angle(u,v)=\cos\angle(Au,Av)=\frac{\sqrt{6}}{6}.\]
\step{Prove the observed behavior for every pair.}
For arbitrary $u,v$, $(Au)^T(Av)=u^TA^TAv=u^Tv$. Taking $v=u$ proves norm preservation, and dividing by nonzero norms proves angle preservation for nonzero vectors.
\answer{$A$ preserves lengths and angles; the numerical experiments are instances of this identity.}

\exercise{6.1.33}{Inner-product and norm calculations}
\step{Compute two actual random examples with one fixed nonzero direction.}
Using seed $61033$ and integer draws from $[-3,3]$, take $v=\left[\begin{array}{r}2\\0\\0\\1\end{array}\right]$, with $v\cdot v=5$. Each scalar below multiplies this same explicitly given column.
\[x=\left[\begin{array}{r}3\\-1\\0\\2\end{array}\right],\quad y=\left[\begin{array}{r}-2\\-3\\1\\1\end{array}\right],\]
\[T(x)=(\frac{8}{5})v,\quad T(y)=(- \frac{3}{5})v,\]
\[T(x+y)=(1)v,\quad T(10x)=(16)v.\]
\[x=\left[\begin{array}{r}-2\\3\\1\\3\end{array}\right],\quad y=\left[\begin{array}{r}-1\\2\\3\\-1\end{array}\right],\]
\[T(x)=(- \frac{1}{5})v,\quad T(y)=(- \frac{3}{5})v,\]
\[T(x+y)=(- \frac{4}{5})v,\quad T(10x)=(-2)v.\]
\step{Verify the conjectured linearity algebraically.}
For scalars $a,b$ and any inputs $x,y$, distributivity gives
\[T(ax+by)=\frac{(ax+by)\cdot v}{v\cdot v}v=aT(x)+bT(y).\]
The fixed denominator is nonzero because $v\ne0$.
\answer{Projection onto a fixed nonzero direction is linear; the displayed experiments satisfy additivity and scalar homogeneity exactly.}

\exercise{6.1.34}{Inner-product and norm calculations}
\step{Read the free variables from an exact row reduction.}
\[\operatorname{rref}(A)=\left[\begin{array}{rrrrr}1&0&5&0&- \frac{1}{3}\\0&1&1&0&- \frac{4}{3}\\0&0&0&1&\frac{1}{3}\\0&0&0&0&0\\0&0&0&0&0\end{array}\right].\]
The pivot columns are $1,2,4$, leaving free coordinates $3,5$. A null-space basis, with the second vector scaled by $3$ to avoid fractions, and a nonzero-row basis for the row space give
\[N=\left[\begin{array}{rr}-5&1\\-1&4\\1&0\\0&-1\\0&3\end{array}\right],\qquad R=\left[\begin{array}{rrrrr}3&0&15&0&-1\\0&3&3&0&-4\\0&0&0&3&1\end{array}\right].\]
The two columns of $N$ have independent free-coordinate choices, and the three rows of $R$ are scaled nonzero rows of RREF, so they are bases of the required spaces.
\step{Multiply the two basis matrices.}
\[RN=\left[\begin{array}{rr}0&0\\0&0\\0&0\end{array}\right].\]
Each entry is the dot product of one row-space basis vector with one null-space basis vector. All are zero, illustrating $(\operatorname{Row}A)^\perp=\Nul A$ from Theorem 3. The additional check $AN=0$ verifies the null-space computation directly.
